\documentclass{article}

\usepackage{fontspec}
\usepackage{gensymb}
\usepackage{graphicx}
\usepackage{hyperref}
\usepackage{cleveref}
\usepackage{titlesec}
\usepackage{titling}
\usepackage{minted}
\usepackage{natbib}
\usepackage{xcolor}
\usepackage{subcaption}
\setlength{\parindent}{0pt}

% Make paragraphs like subsubsubsections
% \setcounter{secnumdepth}{4}
% \titleformat{\paragraph}
% {\normalfont\normalsize\bfseries}{\theparagraph}{1em}{}
% \titlespacing*{\paragraph}
% {0pt}{3.25ex plus 1ex minus .2ex}{1.5ex plus .2ex}

\pretitle{\begin{center} \includegraphics[width=200pt]{shitty-logo.png} \\ \vspace{3em} \LARGE}
\title{Telekinetic Cowboy Country}
\author{
  Emilio Cobos Álvarez \\
  \small \href{mailto:emilio@crisal.io}{emilio@crisal.io}
}

\date{\today}

\begin{document}

\maketitle

\begin{center}
  {\bf Contracts:} UBOSOFT Inc. (not associated in any way with Ubisoft).
\end{center}

\begin{abstract}
  {\it UBOSOFT Inc} is a game publishing company that wants to develop a game
  called {\it Telekinetic Cowboy Country}\footnote{thanks for
  https://videogamena.me/ for the game name.}. \\

  {\it Crisal.io} is a game development contractor company that aims for best
  quality.
\end{abstract}

\clearpage
\tableofcontents

\clearpage

\section{Background}

  Crisal.io is a game development company with a huge reputation in the sector,
  after being responsible of huge successes over the last few years like {\it
  Surprise Magic Mania} (2013) and {\it Inept Wagon Scam} (2015). \\

  UBOSOFT is a game publishing company which is diversifying their efforts and
  are trying to launch a few games on their own. \\

  During the 2016 Game Developers Conference, both companies had a huge
  information exchange about ambitious projects and ideas, which resulted in a
  project called {\it Telekinetic Cowboy Country}\footnote{Oh, wow, you can't
  imagine how much fun I'm having writing this.}.

\section{Objectives and desired outcome}

  The main objective is having the game launched by 2020, and have a pretty good
  sales report. \\

  There are a few stages planned, including having a working a Beta version by
  the end of 2019. \\

  This is an ambitious project and, as such, is also a risky one. We intend to
  innovate in the graphics area, trying to add performant and never-seen-before
  special effects. As such, the plan is also working closely with the Unreal
  Engine developers, contributing features and improvements to their graphics
  pipeline. \\

  Crisal.io depends on their high quality titles to sustain the company, and
  UBOSOFT hasn't have a great last few years, though it's still one of the
  leading game publishing companies. As such, is in the best interest of both
  companies to have a greatly successful title as the outcome of this project.

\section{Scope and restrictions of the project}

  As discussed before, both companies want a hugely successful title and have a
  good amount of resources and experience to achieve so. \\

  Crisal.io is an experienced company in game development, though it doesn't
  have a huge amount of designers. As such, it may be possible to outsource
  character and some episode design to third parties. \\

  Also, the game wants to innovate in the space of graphics and performance, but
  it's definitely not in the scope to develop a new rendering engine from
  scratch. The Crisal developers are experienced graphics programmers, and avid
  contributors to the Unreal Engine project, which we plan to use for this
  project. \\

  We intend to publish the game for PC, Linux and Mac during the initial launch,
  assuming modern graphics pipelines (Vulkan and Metal-based). \\

  After that initial launch, we may be able to look into supporting older
  Graphics card via OpenGL/DirectX, or video game consoles, depending on market
  analysis as the date approaches.

\section{Project tolerances}

  The project scope is not really discutible. The end result must be a
  production-quality game. We may consider making the scope broader and
  prolonging the business relationships among parties with DLCs and extra
  content. \\

  The schedule is probably somewhat flexible, but we've scheduled mid-2020 as
  the latest day possible for the first release, that is, about 20\% schedule
  tolerance. \\

  Regarding cost tolerance, the cost is still being discussed, and may vary a
  lot on later stages of the project. Still the budget is initially \$400.000,
  and never going above \$500.000, at least initially.

\section{Risks}

  The project is fairly risky. We intend to develop technology on the bleeding
  edge, and targets may move. \\

  The game concept is also unusual, so even if we succeed in the technical
  aspects, there are a lot of different errors that may happen and render our
  project useless, like market disliking the innovation in the game concept
  area. \\

  Still, both companies are able and willing to go ahead with this project, and
  wish that their vision comes true. \\

  In any case, this is still not the end for any of the companies. On one hand,
  UBOSOFT is one of the biggest game publishers, and have a bunch of other
  projects at the same time. Crisal would still be one of the leader consultancy
  companies, though its reputation may be severely damaged if the project is
  unsuccessful.

\section{Other interested parties}

  Obviously, developers and publishers are the most interested parties inside
  this project. \\

  We've received some interest from graphic card manufacturers in order to help
  them test and profile their new Metal/Vulkan graphics drivers. \\

  Also, a few gaming communities have expressed interest after the announcement
  of the project.

\section{Business case}

  As discussed, both companies are doing a quite risky investment here. UBOSOFT
  expects this title to be a hit and obtain a fair amount of revenue. \\

  Crisal would only get the contract plus the reputation if the title is
  successful. Also, it's possible to arrange follow-up extra-content contracts
  if the game succeeds, in order to offer DLCs and similar stuff. \\

  The business model in the Game sector is far from innovative. The game will be
  launched over Steam for PC, Linux and OSX platforms, and on the usual game
  stores for consoles. \\

  We'll consider making extra content available via DLCs as we see fit after the
  initial release. \\

\end{document}
